\documentclass[11pt]{article}
\usepackage[margin=0.8in]{geometry}
\usepackage[T1]{fontenc}
\usepackage{lmodern,amsmath,amssymb,amsthm,booktabs,array}
\usepackage[colorlinks=true,linkcolor=blue,urlcolor=blue,citecolor=blue]{hyperref}
\hypersetup{pdftitle={Binary columns, polynomial reductions, and quartic harmonic means},pdfauthor={Blake Stone},pdfsubject={Four recorded source assertions in three families, with complete Lean proofs}}
\setlength{\parindent}{0pt}
\setlength{\parskip}{0.38em}
\setlength{\emergencystretch}{1em}
\newtheorem{theorem}{Theorem}[section]
\newtheorem{lemma}[theorem]{Lemma}
\newtheorem{proposition}[theorem]{Proposition}
\newtheorem{corollary}[theorem]{Corollary}
\newcommand{\N}{\mathbb N}
\newcommand{\Z}{\mathbb Z}
\newcommand{\Q}{\mathbb Q}
\newcommand{\R}{\mathbb R}
\begin{document}
\begin{center}
{\LARGE\bfseries Binary columns, polynomial reductions,\\[3pt]
and quartic harmonic means}\par
\vspace{0.5em}
{\large Four source assertions in three families: complete Lean proofs}\par
\vspace{0.5em}
{\large Blake Stone}\par
Proofs developed with ChatGPT\par
November 27, 2026
\end{center}

\textbf{Abstract.}
We prove four specific conjectural assertions in three mathematical families.
For the increasing array of positive integers with a fixed number of binary
ones, we prove Kimberling's eventual column-doubling conjecture and the exact
least starting row $\max(1,j-1)$ for every column $j$. For the two components
of central-binomial prefix polynomials reduced by $X^2\mapsto X+1$, a direct
divisibility certificate proves both of Mathar's recorded recurrences.
Finally, exact rational moment inequalities establish the generating function
of the floors of harmonic means between consecutive fourth powers. The last
proof supplies the corrected same-index floor formula, the valid recurrence
cutoff, and a full source-to-power-series bridge.
Each result is verified from its literal source definition in Lean, with
separate fresh author and reviewer builds, exact source snapshots, and a
mathematical review outside its authoring team. Prior methods and related
results are credited. Bounded searches do not establish worldwide priority.

\begin{center}
\begin{tabular}{@{}l>{\raggedright\arraybackslash}p{4.3in}@{}}
\toprule
Source & Exact scope\\
\midrule
A067576 & Eventual doubling in every binary-weight column, with the exact least starting row.\\
A192250/A192251 & Two component recurrences from the literal terminating polynomial-reduction algorithm.\\
A227013 & The conjectured harmonic-mean generating function, with corrected closed-form and recurrence indices.\\
\bottomrule
\end{tabular}
\end{center}

The two component recurrences share one proof and count as one family.
The onset formula, intermediate inequalities, recurrence consequences, and
indexing corrections are not additional problem counts. These are proofs
of precise recorded assertions; the established theories used to prove them
are not claimed as new.

\section{Eventual doubling in the binary-weight array}

For positive integers $i,j$, let $T(i,j)$ be the $j$th positive integer
whose ordinary binary expansion contains exactly $i$ ones, listed in
increasing order. This is the array of OEIS A067576, introduced by
Wilson in 2002. Kimberling conjectured in November 2025 that every
column of $T(i,j)+1$ eventually forms a geometric progression of ratio
two \cite{oeis067576}. We prove the conjecture and determine its exact
starting row.

\begin{theorem}\label{a067576:main}
For positive $i,j$ with $i\ge j-1$,
\begin{equation}\label{a067576:formula}
 T(i,j)+1=(2^j-1)2^{i+1-j}.
\end{equation}
Consequently the least positive row $r$ such that
\[
 T(i+1,j)+1=2\bigl(T(i,j)+1\bigr)\qquad\text{for every }i\ge r
\]
is exactly $\max(1,j-1)$.
\end{theorem}

\subsection*{Exact ranks from binary prefix counts}
Let $w(n)$ be the number of ones in the binary expansion of the natural
number $n$, including $w(0)=0$, and put
\[
 C_i(m)=\#\{n\in\mathbb N: n<m,\ w(n)=i\}.
\]
The identities $w(2n)=w(n)$ and $w(2n+1)=w(n)+1$ imply
\begin{align}
 C_{i+1}(2m)&=C_{i+1}(m)+C_i(m),\label{a067576:even}\\
 C_{i+1}(2m+1)&=C_{i+1}(m+1)+C_i(m).\label{a067576:odd}
\end{align}
Indeed, the even integers below $2m$ correspond to integers below $m$
of unchanged weight; the odd integers correspond to integers below $m$
of weight one less. Below $2m+1$, the even part additionally contains
$2m$, giving the second equation.

Binary induction proves $2^{w(n)}\le n+1$. At an even successor step,
$2^{w(2n)}=2^{w(n)}\le n+1\le2n+1$; at an odd step,
$2^{w(2n+1)}=2\cdot2^{w(n)}\le2n+2$. The base $n=0$ is equality.
Also $w(2^i-1)=i$, since the binary expansion is a string of $i$ ones.
It follows that
\begin{equation}\label{a067576:minimum}
 C_i(m)=0\quad(m<2^i),\qquad C_i(2^i)=1.
\end{equation}
For the first equality, every $n<m$ satisfies $n+1\le m<2^i$ and
therefore cannot have weight $i$. The second follows by adjoining the
single integer $2^i-1$.

Induction in \eqref{a067576:even} now gives
\begin{equation}\label{a067576:nearfull}
 C_i(2^{i+1})=i+1.
\end{equation}
The case $i=0$ counts only zero. For the induction step, the two terms
in \eqref{a067576:even} are $C_{i+1}(2^{i+1})=1$ and
$C_i(2^{i+1})=i+1$. The two largest integers below $2^{i+1}$ have
weights
\[
 w(2^{i+1}-2)=i,\qquad w(2^{i+1}-1)=i+1.
\]
Subtracting their contributions from \eqref{a067576:nearfull} shows
\begin{equation}\label{a067576:base}
 w(2^{i+1}-2)=i,\qquad C_i(2^{i+1}-2)=i.
\end{equation}

\begin{lemma}\label{a067576:rank}
For all $k,d\ge0$, set $v(k,d)=(2^{k+1}-1)2^d-1$. Then
\[
 w(v(k,d))=k+d,\qquad C_{k+d}(v(k,d))=k.
\]
\end{lemma}
\begin{proof}
For $d=0$, this is \eqref{a067576:base}. Suppose it holds for $d$,
and write $v=v(k,d)$. Then
\[
 v(k,d+1)=2v+1,\qquad
 v+1=(2^{k+1}-1)2^d<2^{k+d+1}.
\]
The weight of $2v+1$ is $k+d+1$. Equations
\eqref{a067576:odd} and \eqref{a067576:minimum} give
\[
 C_{k+d+1}(2v+1)
 =C_{k+d+1}(v+1)+C_{k+d}(v)=0+k.
\]
This proves the induction step.
\end{proof}

For every positive $i$, the integers $(2^i-1)2^a$, as $a$ varies, show
that the weight-$i$ row is infinite. Its increasing enumeration is
therefore characterized by membership and the exact number of smaller
members. Lemma~\ref{a067576:rank} says that $v(k,d)$ is the $(k+1)$st
integer of weight $k+d$. For $j\ge1$ and $i\ge j-1$, take
$k=j-1$ and $d=i+1-j$. Then $k+d=i$, and the lemma gives
\eqref{a067576:formula}. Multiplying its right side by two gives the
formula at row $i+1$, proving eventual doubling from
$\max(1,j-1)$ onward.

\subsection*{The starting row is minimal}
For $k\ge0$, consider
\[
 u=2^{k+2}+(2^k-1).
\]
Adding the isolated leading bit shows $w(u)=k+1$. More generally, for
$0\le n<2^r$ one has $w(2^r+n)=1+w(n)$, and hence for
$0\le m\le2^r$,
\begin{equation}\label{a067576:leading}
 C_{i+1}(2^r+m)=C_{i+1}(2^r)+C_i(m).
\end{equation}
The equality splits the count into the integers below $2^r$ and those
having the new leading bit. Applying it with $i=k$, $r=k+2$ and
$m=2^k-1$, and then using \eqref{a067576:nearfull} and
\eqref{a067576:minimum}, yields
\[
 C_{k+1}(u)=C_{k+1}(2^{k+2})+C_k(2^k-1)=k+2.
\]
Thus $u=T(k+1,k+3)$, so
\[
 T(k+1,k+3)+1=5\cdot2^k.
\]
At the next row, the already proved formula gives
\[
 T(k+2,k+3)+1=2^{k+3}-1=8\cdot2^k-1
 \ne10\cdot2^k=2\bigl(T(k+1,k+3)+1\bigr).
\]
For a column $j\ge3$, choosing $k=j-3$ proves failure at row $j-2$.
Any earlier proposed onset would include this transition and therefore
cannot work. For $j=1,2$, doubling already holds from the least positive
row, namely one. This completes Theorem~\ref{a067576:main}.

\paragraph{Formalization and attribution.}
The formal weight is \texttt{Nat.bitIndices.length}, and the source
entry is the zero-based selector \texttt{Nat.nth} at index $j-1$.
The proof establishes infinite positive rows, strict increase, exact
membership, and exact rank, as well as the universal formula and the
\texttt{IsLeast} onset statement. All four source modules compile with
Lean 4.19.0 and pinned mathlib, and their principal theorem audits use
only the standard axioms \texttt{propext}, \texttt{Classical.choice},
and \texttt{Quot.sound}.

Binary-weight ranking is classical. A036563 records the diagonal
special case, and A066884 presents this array with reversed
antidiagonals. The paper of Dobri\'c, Skyers and Stanley develops
polynomial-time weight ranking and unranking \cite{dobric2016}.
The minimum-onset formula was also already posed as an admitted Lean
statement in open-math-agents \cite{oma067576}; that file contains no
proof. We claim a complete proof of the recorded eventual-doubling
assertion and its exact onset, with no claim of first formulation or
global historical priority. The package counts as one source-conjecture
resolution.


\section{Two central-binomial polynomial-reduction recurrences}
\label{sec:a192250}

Let $c_k=\binom{2k}{k}$ and, for $n\geq 1$, set
\[
 P_n(X)=\sum_{k=0}^{n-1}c_kX^k.
\]
The polynomial reduction used by Kimberling in OEIS A192232 repeatedly
writes a polynomial as $X^2t(X)+r(X)$, where $\deg r<2$, and replaces it by
$(X+1)t(X)+r(X)$ until a polynomial of degree at most one is reached
\cite{OEISA192232}. Write the result of reducing $P_n$ as
\[
 R_n(X)=a_n+b_nX.
\]
Then $a_n$ is A192250 and $b_n$ is A192251
\cite{OEISA192250,OEISA192251}. The original offsets are both~1.
In particular, the source program makes its polynomial list starting at
degree~0, and reads its first element at sequence index~1. Thus the input
at index $n$ has $n$ terms, ending in $c_{n-1}X^{n-1}$; it is not $P_{n+1}$.
The first five pairs $(a_n,b_n)$ are
\[
(1,0),\quad(1,2),\quad(7,8),\quad(27,48),\quad(167,258).
\]

\begin{theorem}
For either $y_n=a_n$ or $y_n=b_n$, and every $n\geq4$,
\begin{equation}
\label{eq:a192250-mathar}
\begin{split}
 &(n-1)(n-2)y_n-(5n-7)(n-2)y_{n-1}\\
 &\qquad{}-2(2n-3)(3n-8)y_{n-2}
       +4(2n-3)(2n-5)y_{n-3}=0.
\end{split}
\end{equation}
\end{theorem}

These are the two identical recurrences recorded as conjectures by
R.~J.~Mathar on May~4, 2014. The range $n\geq4$ is the full range in
which the recurrence uses only terms at the source's indices $1,2,\ldots$.
The two assertions belong to one polynomial family.

\begin{proof}
Put $M(X)=X^2-X-1$. First, each source step satisfies
\[
 X^2t+r-\bigl((X+1)t+r\bigr)=M(X)t.
\]
It therefore preserves the remainder modulo $M$. If the input degree is
at least two, the new degree is at most the old degree minus one:
$\deg t=\deg p-2$, while $\deg r<2$ and
$\deg((X+1)t)\leq\deg t+1$. Hence the algorithm terminates.
Its terminal polynomial has degree less than two and is congruent to
the input modulo the monic quadratic $M$, so uniqueness of polynomial
division identifies it with the remainder modulo $M$. This supporting
observation is standard; it is also Lemma~1 in Wiktorowicz's 2026 note
\cite{Wiktorowicz2026Reduction}. We include its proof here to specify the
source interpretation completely.

Fix $n\geq4$, and abbreviate
\[
 A=(n-1)(n-2),\qquad B=(5n-7)(n-2),
\]
\[
 C=2(2n-3)(3n-8),\qquad D=4(2n-3)(2n-5).
\]
We claim the following exact identity in $\mathbb Z[X]$:
\begin{equation}
\label{eq:a192250-certificate}
 A P_n-B P_{n-1}-C P_{n-2}+D P_{n-3}
   =D c_{n-3}X^{n-3}(X^2-X-1).
\end{equation}
To verify it, first observe
\[
 A-B-C+D=0,\qquad A-B=-2(2n-3)(n-2).
\]
The elementary central-binomial identity
\[
 (k+1)c_{k+1}=2(2k+1)c_k
\]
follows directly from $c_k=(2k)!/(k!)^2$. Applying it successively at
$k=n-3$ and $k=n-2$ gives
\[
 A c_{n-1}=D c_{n-3},\qquad
 (A-B)c_{n-2}=-D c_{n-3}.
\]
Expanding the three last terms of $P_n$ around $P_{n-3}$ now yields
\begin{align*}
 A P_n-B P_{n-1}-C P_{n-2}+D P_{n-3}
 &= (A-B-C+D)P_{n-3}\\
 &\quad +(A-B-C)c_{n-3}X^{n-3}\\
 &\quad +(A-B)c_{n-2}X^{n-2}
       +A c_{n-1}X^{n-1}\\
 &=D c_{n-3}\bigl(-X^{n-3}-X^{n-2}+X^{n-1}\bigr),
\end{align*}
which is exactly \eqref{eq:a192250-certificate}.

Remainder modulo a monic polynomial is $\mathbb Z$-linear.
Taking remainders in \eqref{eq:a192250-certificate}, and using the source
identification established above, gives
\[
 A R_n-B R_{n-1}-C R_{n-2}+D R_{n-3}=0.
\]
Taking the constant coefficient proves \eqref{eq:a192250-mathar} for
$a_n$, and taking the coefficient of $X$ proves it for $b_n$.
\end{proof}

\paragraph{Formal verification.}
The Lean definition \path{sourceReduction} implements the actual
division-and-substitution procedure by well-founded recursion on the
polynomial's natural degree. The theorem
\path{sourceStep_natDegree_lt} proves the decrease needed for every
recursive call, and \path{sourceReduction_eq_remainder} proves its
equivalence to the monic remainder. The complete polynomial certificate
is proved in \path{polynomial_certificate}. The endpoints
\path{a192250_conjecture} and \path{a192251_conjecture} use the
literal algorithm and quantify over every $n\geq4$ in the original
indexing. All three local modules are explicit default build roots.
A fresh source-only build with Lean~4.19.0 and the pinned mathlib revision
succeeded; all audited endpoints use only the standard axioms
\texttt{propext}, \texttt{Classical.choice}, and \texttt{Quot.sound}.

\paragraph{Scope and prior work.}
The general reduction lemma, the Fibonacci description of the reduced
monomials, and Campbell's 2016 sum formula for A192251 are supporting
prior material, not additional new results. Wiktorowicz's explicit
theorem and table concern polynomial input sequences and
constant-coefficient recurrences. Here the input is the central-binomial
sequence, and the selected endpoints are Mathar's two specific
polynomial-coefficient recurrences. The bounded source and semantic
search is recorded in the package; no claim of worldwide priority is
made.


\section{A227013: harmonic means between fourth powers}
\label{sec:quartic-harmonic}

For $n\ge1$, define the literal harmonic mean and its floor by
\[
 d_n=n^4-(n-1)^4,\qquad
 H_n=\sum_{k=(n-1)^4+1}^{n^4}\frac1k,\qquad
 M_n=\frac{d_n}{H_n},\qquad a_n=\lfloor M_n\rfloor.
\]
Kimberling's A227013 records a conjectural rational generating function,
closed form, and recurrence \cite{oeis227013}. We prove the exact generating
function, while explicitly correcting the index conventions of the other
formulas. This is one source family.

\begin{theorem}\label{thm:quartic-main}
Put
\[
 Q(n)=n^4-2n^3+\frac53n^2-\frac23n
      =\frac{n(n-1)(3n^2-3n+2)}3.
\]
Then $a_1=1$, $a_2=6$, and $a_n=\lfloor Q(n)\rfloor$ for every $n\ge3$.
The exact formal generating function, with coefficient $x^j$ equal to
$a_{j+1}$, is
\begin{equation}\label{eq:quartic-gf}
 \sum_{j\ge0}a_{j+1}x^j=
 \frac{-1-2x-22x^2-23x^3-20x^4-4x^5+2x^6-3x^7+x^8}
 {(x-1)^5(1+x+x^2)}.
\end{equation}
For every $n\ge10$,
\begin{equation}\label{eq:quartic-recurrence}
 a_n=4a_{n-1}-6a_{n-2}+5a_{n-3}-5a_{n-4}
       +6a_{n-5}-4a_{n-6}+a_{n-7}.
\end{equation}
\end{theorem}

\subsection{A finite rational moment bound}

Consider any consecutive block $a+1,\ldots,a+d$ with $a\ge0$ and
positive integer $d$. Set
\[
 c=a+\frac{d+1}2,\qquad r=\frac{d-1}2,\qquad t_j=j-r\quad(0\le j<d).
\]
The reflection $j\mapsto d-1-j$ negates $t_j$. Thus
\begin{equation}\label{eq:quartic-pairing}
 \sum_{j=0}^{d-1}\frac1{a+1+j}
 =\sum_{j=0}^{d-1} B(c,t_j),\qquad
 B(c,t)=\frac12\left(\frac1{c+t}+\frac1{c-t}\right).
\end{equation}
Here $c>r\ge|t_j|$, so all denominators are positive.
The exact identities
\begin{align}
 B(c,t)-\left(\frac1c+\frac{t^2}{c^3}+\frac{t^4}{c^5}+\frac{t^6}{c^7}\right)
 &=\frac{t^8}{c^7(c^2-t^2)}\ge0,\label{eq:quartic-lower-pair}\\
 \left(\frac1c+\frac{t^2}{c^3}+\frac{t^4}{c^3(c^2-r^2)}\right)-B(c,t)
 &=\frac{t^4(r^2-t^2)}{c^3(c^2-r^2)(c^2-t^2)}\ge0
 \label{eq:quartic-upper-pair}
\end{align}
hold for every $|t|\le r$. No infinite expansion or analytic remainder
estimate is involved.

The averaged central moments of $0,\ldots,d-1$ are
\begin{align*}
 v_2&=\frac1d\sum_{j<d}t_j^2=\frac{d^2-1}{12},\\
 v_4&=\frac1d\sum_{j<d}t_j^4=\frac{(d^2-1)(3d^2-7)}{240},\\
 v_6&=\frac1d\sum_{j<d}t_j^6=\frac{(d^2-1)(3d^4-18d^2+31)}{1344}.
\end{align*}
These identities follow by expanding the powers and summing the ordinary
powers through degree six. In the formal proof each ordinary power sum
is proved by induction on $d$, and the three displayed identities follow
by polynomial algebra. Summing \eqref{eq:quartic-lower-pair} and
\eqref{eq:quartic-upper-pair} gives the exact bounds
\begin{equation}\label{eq:quartic-LU}
 dL\le H\le dU,\qquad
 L=\frac1c+\frac{v_2}{c^3}+\frac{v_4}{c^5}+\frac{v_6}{c^7},\qquad
 U=\frac1c+\frac{v_2}{c^3}+\frac{v_4}{c^3(c^2-r^2)}.
\end{equation}

\subsection{Uniform separation of the harmonic mean from its floor}

For the quartic block, take $a=(n-1)^4$ and $d=d_n$. Then
\[
 c=(n^2-n+1)^2,\qquad r=(d_n-1)/2,\qquad
 c^2-r^2=n^4((n-1)^4+1).
\]
For $n=t+5$, direct clearing of denominators in \eqref{eq:quartic-LU}
gives the following two exact rational identities:
\begin{align}
 (Q(n)+2/3)L-1
 &=\frac{P(t)}{315(t^2+9t+21)^{14}},\label{eq:quartic-P}\\
 1-Q(n)U
 &=\frac{R(t)}{45(t+5)^2(t^2+9t+21)^6((t+4)^4+1)}.
 \label{eq:quartic-R}
\end{align}
All coefficients of $P$ and $R$, including their constant terms, are
positive. The complete coefficient lists appear in Table~\ref{tab:quartic-cert};
they define the polynomials by $P(t)=\sum_kP_kt^k$ and
$R(t)=\sum_kR_kt^k$. Consequently both right sides are strictly positive
for every $t\ge0$. The identities themselves are checked by polynomial
normalization in Lean, so the table is an exact algebraic certificate.

Since $Q(n)>0$, $d_n>0$, and $H_n>0$, equations
\eqref{eq:quartic-LU}--\eqref{eq:quartic-R} imply
\[
 Q(n)H_n\le Q(n)d_nU<d_n,
 \qquad
 (Q(n)+2/3)H_n\ge (Q(n)+2/3)d_nL>d_n.
\]
Therefore, for every $n\ge5$,
\begin{equation}\label{eq:quartic-sandwich}
 Q(n)<M_n<Q(n)+\frac23.
\end{equation}

For an integer $n$, the fractional part of $Q(n)$ is $1/3$ when
$n\equiv2\pmod3$ and is zero otherwise. This follows by substituting
$n=3k,3k+1,3k+2$ into its polynomial. Thus
\[
 Q(n)+2/3\le\lfloor Q(n)\rfloor+1.
\]
Together with \eqref{eq:quartic-sandwich}, this places $M_n$ in the
same unit interval as $Q(n)$ and proves $a_n=\lfloor Q(n)\rfloor$
for every $n\ge5$.

The remaining literal harmonic means satisfy the following finite rational
comparisons:
\[
\begin{array}{c|r|r|c}
n&d_n&a_n&\text{exact comparison for }H_n\\\hline
1&1&1&H_1=1\\
2&15&6&15/7<H_2\le15/6\\
3&65&40&65/41<H_3\le65/40\\
4&175&152&175/153<H_4\le175/152
\end{array}
\]
They can be checked by multiplying through by the common denominator
of the finitely many positive reciprocals. In Lean all four values
are explicit kernel computations on exact rational numbers. In particular
$\lfloor Q(3)\rfloor=40$ and $\lfloor Q(4)\rfloor=152$, completing
the floor formula at every claimed index.

\begin{table}[p]
\centering
\small
\caption{Complete positive-coefficient certificates for
\eqref{eq:quartic-P} and \eqref{eq:quartic-R}. A dash means coefficient zero.}
\label{tab:quartic-cert}
\begin{tabular}{r|r|r}
$k$&$P_k$&$R_k$\\\hline
24&28&---\\
23&3024&---\\
22&156572&---\\
21&5172156&---\\
20&122369371&---\\
19&2206424754&---\\
18&31496055869&---\\
17&364973853078&---\\
16&3492525189361&---\\
15&27925808790492&---\\
14&188059104987530&26\\
13&1071825627742992&1612\\
12&5182093657336537&46362\\
11&21253816999367298&819836\\
10&73776484843852949&9958385\\
9&215711119121472558&87898188\\
8&527200074613081231&581403224\\
7&1064965904311292064&2927810658\\
6&1749657582705865988&11280438711\\
5&2284392530671416108&33094892742\\
4&2290486550158893280&72791311130\\
3&1671930872839271052&116416933660\\
2&810572143783759280&128020606587\\
1&216472755572592948&86680806370\\
0&18186223121964135&27278343429\\
\end{tabular}
\end{table}

\subsection{The recurrence and exact generating function}

The sequence $b_n=\lfloor Q(n)\rfloor$ is a quartic polynomial plus
a correction of period three:
\[
 b_n=Q(n)-\frac13\mathbf1_{n\equiv2\ (\mathrm{mod}\ 3)}.
\]
If $E$ is the forward shift, the polynomial part is annihilated by
$(E-1)^5$ and the periodic part by $E^3-1$. Both are therefore
annihilated by
\[
 (E-1)^5(E^2+E+1)
 =E^7-4E^6+6E^5-5E^4+5E^3-6E^2+4E-1.
\]
This gives \eqref{eq:quartic-recurrence} whenever all its indices are
at least three, that is, $n\ge10$. The Lean proof expands the three
residue-class polynomials and checks this identity for arbitrary
quotient $k$, then transports it to the literal source sequence.

The first nine terms are
\[
 (a_1,\ldots,a_9)=(1,6,40,152,413,920,1792,3173,5232).
\]
Let $A(x)=\sum_{j\ge0}a_{j+1}x^j$ and
$D(x)=(1-x)^5(1+x+x^2)$.
The recurrence makes every coefficient of $D(x)A(x)$ from degree nine
onward zero; its first nine coefficients are
\[
 (1,2,22,23,20,4,-2,3,-1).
\]
Thus
\[
 D(x)A(x)=1+2x+22x^2+23x^3+20x^4+4x^5-2x^6+3x^7-x^8.
\]
Negating numerator and denominator gives exactly
\eqref{eq:quartic-gf}. The denominator has constant coefficient $-1$,
so it is a unit in $\mathbb Z[[x]]$ and determines the quotient uniquely.
This proves Theorem~\ref{thm:quartic-main}.

\paragraph{Source indexing.}
The source's generating function has constant coefficient $1$ despite
its offset of one; our convention is explicitly $\sum_{j\ge0}a_{j+1}x^j$.
The separately printed expression
\[
 \frac{47}{9}+14n+\frac{41n^2}{3}+6n^3+n^4
       -\frac29\cos(2\pi n/3)
\]
equals $a_{n+2}$ for positive integer $n$. This follows by substituting
$n+2$ into $Q$ and using the cosine's residue-class values $1,-1/2,-1/2$.
At $n=1$ it gives $40$, whereas $a_1=1$. It must not be interpreted
as a correct same-index formula. The raw recurrence cutoff ``$n>2$''
also fails to specify valid previous positive indices, and even the
well-defined case $n=9$ fails: its right side is $5233$ instead of
$a_9=5232$. The corrected unbounded cutoff is $n\ge10$.
The exact generating-function assertion is the selected resolution;
these index clarifications are not additional problem counts.

\paragraph{Formal scope and prior work.}
Eight modules define the finite reciprocal sum and harmonic mean,
prove the central moments and rational bounds, check the polynomial
certificates, identify the floor at every positive index, and establish
the full formal generating function. The endpoint
\path{QuarticHarmonic.full_resolution} includes the literal coefficient
identification and the exact source numerator--denominator identity.
All principal dependencies are the three standard Lean logical axioms;
there is no admitted inequality, asymptotic hypothesis, or floating-point
calculation in the proof.

The nearby cubic problem A227012 and a general theorem for cubic block
functions already have an exact proof and Lean development in
\path{rajdeepbaror/kimberling-a227012-lean} \cite{prior-cubic-harmonic}.
We inspected that work and excluded the cubic problem from this batch.
Its hypotheses do not cover the present fourth-power blocks.
Moment identities and symmetric reciprocal bounds are elementary
classical tools. Existing jOEIS and LODA code already encodes the proposed
recurrence; the contribution here is its complete proof from the literal
harmonic mean. The bounded exact and semantic prior-work screen is
recorded separately and does not establish worldwide priority.

\clearpage
\section{Formal verification and provenance}

The collection contains three independently buildable projects with fifteen
local proof modules. Every module is an explicit default Lake root. The
compiler is Lean 4.19.0 and every package pins mathlib commit
\begin{center}
\texttt{c44e0c8ee63ca166450922a373c7409c5d26b00b}.
\end{center}

\begin{center}
\begin{tabular}{@{}lcp{3.8in}@{}}
\toprule
Package & Modules & Complete formal conclusion\\
\midrule
A067576 & 4 & Literal positive increasing binary-weight enumeration, exact formula, eventual doubling, and the least onset for every column.\\
A192250 & 3 & Termination and correctness of the literal source substitution, followed by both recurrences at every original index $n\ge4$.\\
A227013 & 8 & Literal finite harmonic mean, exact rational sandwich, all positive floor values, and the full source generating function.\\
\bottomrule
\end{tabular}
\end{center}

Each package was built twice from source and pinned configuration in separate
fresh directories, without copying its local compiled proof objects.
Only pinned external dependency caches were shared. Build logs record all
source hashes and axiom reports. All principal endpoints depend only on
\path{propext}, \path{Classical.choice}, and \path{Quot.sound}. No theorem
is admitted, no custom axiom is introduced, and no proof invokes
\path{native_decide}. The four small harmonic-mean values are computed
by Lean's kernel using exact rational arithmetic.

A reviewer who did not contribute to a package's proof sources checked its
mathematics, exact source statements, indices, and endpoint scope, and ran
its separate clean build. The reviews also repeated literal enumeration or
substitution computations and checked the quartic polynomial certificates
independently. These are independent reviews within this collaboration,
not journal peer review. The finite checks detect transcription or source
interpretation errors; the universal conclusions are proved mathematically
and checked by Lean.

\paragraph{Known work and result counting.}
Binary ranking and unranking are classical, and the exact minimum-onset
formula was already formulated as an unproved Lean target. The polynomial
reduction-to-remainder observation and known Fibonacci-weighted descriptions
are credited, including Wiktorowicz's reduction lemma and Campbell's formula.
The prior complete proof of A227012 and the general cubic theorem were found
during screening and excluded; their hypotheses do not cover the quartic
blocks treated here. Existing source formulas and numerical generators are
not presented as new discoveries.

This batch covers four OEIS entries in three related families. It resolves
the selected assertions, not every nearby question or an arbitrary-parameter
extension. In particular, the quartic result is a complete proof of one
specific harmonic-mean generating function; it is not a proof of the broad
recurrence question for all polynomial block functions.

\paragraph{Source record.}
The packages include immutable official OEIS exports and hashes, exact and
semantic searches through both web engines, global GitHub searches, and
checks of current formal-proof collections. Relevant papers and code were
inspected for actual hypotheses and conclusions. No earlier proof of the
selected exact assertions was found in those bounded searches. Search caps,
indexing gaps, and inaccessible material limit historical coverage; the
records do not establish worldwide priority. No OEIS edits or external
announcements are part of the repository publication.

\begin{thebibliography}{12}
\raggedright
\bibitem{oeis067576}
R. G. Wilson v; conjecture by C. Kimberling.
OEIS A067576, revision 57, November 11, 2025.
\url{https://oeis.org/A067576}.

\bibitem{dobric2016}
V. Dobri\'c, M. Skyers, and L. J. Stanley.
\emph{Polynomial Time Computable Triangular Arrays For Almost Sure Convergence}.
arXiv:1603.04896, 2016.
\url{https://arxiv.org/abs/1603.04896}.

\bibitem{oma067576}
GuillaumeBAUD-coder, open-math-agents.
\texttt{statements/oeis-a067576.lean}, admitted conjecture statement,
observed commit \texttt{3a0039bdf1da87d3915b73dabafd68415871f2e4}.
\url{https://github.com/GuillaumeBAUD-coder/open-math-agents}.

\bibitem{OEISA192250}
The OEIS Foundation Inc., A192250,
\emph{The On-Line Encyclopedia of Integer Sequences},
\url{https://oeis.org/A192250}, accessed November 27, 2026.

\bibitem{OEISA192251}
The OEIS Foundation Inc., A192251,
\emph{The On-Line Encyclopedia of Integer Sequences},
\url{https://oeis.org/A192251}, accessed November 27, 2026.

\bibitem{OEISA192232}
The OEIS Foundation Inc., A192232,
\emph{The On-Line Encyclopedia of Integer Sequences},
\url{https://oeis.org/A192232}, accessed November 27, 2026.

\bibitem{Wiktorowicz2026Reduction}
R. Wiktorowicz,
\emph{Generating functions for reductions by $x^2\to x+1$: proofs of the empirical formulas},
2026, \url{https://oeis.org/A192232/a192232.pdf}.

\bibitem{oeis227013} C.~Kimberling. OEIS A227013, revision 21, November 3, 2024.
\url{https://oeis.org/A227013}.
\bibitem{prior-cubic-harmonic} \path{rajdeepbaror/kimberling-a227012-lean}. \emph{Clark Kimberling's Harmonic-Mean Conjectures (OEIS A227012 and General Integer Cubics)},
repository inspected at commit \texttt{f33c98bac3c0eabccf7c66e339d661494960c523}.
\url{https://github.com/rajdeepbaror/kimberling-a227012-lean}.

\end{thebibliography}
\end{document}
